\chapter{Partitioning, Bucketing, and Pruning}
\label{ch:pruning}

\begin{chapterintro}
The cheapest data to move is data you never read. Partition pruning skips whole
directories; dynamic partition pruning does it using values discovered from the
other side of a join; bucketing lays data out so a join needs no shuffle at all.
Each of these is easy to defeat --- a function on the partition column, a
\texttt{CAST}, a stray \texttt{SELECT *} --- and the plan tells you when you
have. This chapter covers each mechanism and how to confirm it is working.
\end{chapterintro}

\section{Partition pruning and the pushdown killers}
\tbd

\section{Dynamic partition pruning}
\tbd

\section{Bucketing to avoid a shuffle}
\tbd

\section{\texttt{repartition} versus \texttt{coalesce}}
\label{sec:repartition-coalesce}
\tbd

\section{Column pruning and \texttt{ReadSchema}}
\tbd

\section{Summary}
\tbd

\exercisehead
\begin{exercises}
\item Write the same filter three ways --- a function on the partition column, a
      \texttt{CAST}, and a plain literal --- and check which keeps
      \texttt{PartitionFilters} populated in the scan node.
\end{exercises}
